% ============================================================
% Standalone manuscript
% Compiler: pdfLaTeX
% ============================================================

\documentclass[12pt]{article}

\usepackage[T1]{fontenc}
\usepackage[utf8]{inputenc}
\usepackage{lmodern}
\usepackage{amsmath,amsthm,amsfonts,amssymb}
\usepackage[a4paper,margin=27mm]{geometry}
\usepackage{microtype}
\usepackage[hidelinks]{hyperref}

\numberwithin{equation}{section}

\newtheorem{theorem}{Theorem}[section]
\newtheorem{lemma}[theorem]{Lemma}
\newtheorem{proposition}[theorem]{Proposition}
\newtheorem{corollary}[theorem]{Corollary}

\theoremstyle{definition}
\newtheorem{definition}[theorem]{Definition}
\newtheorem{example}[theorem]{Example}
\theoremstyle{remark}
\newtheorem{remark}[theorem]{Remark}

\newcommand{\Z}{\mathbb Z}
\newcommand{\one}[1]{\mathbf 1_{\{#1\}}}
\newcommand{\Ccal}{\mathcal C}
\newcommand{\Rcal}{\mathcal R}
\newcommand{\Zcal}{\mathcal Z}
\newcommand{\Ecal}{\mathcal E}
\newcommand{\Bcal}{\mathcal B}

\hypersetup{
  pdftitle={
    Exact Formulas for Restricted Coprime Representations
    of Even Integers with Squarefree Modulus 6P
  },
  pdfauthor={Andres M. Salazar},
  pdfkeywords={
    additive representations,
    coprimality,
    Chinese remainder theorem,
    canonical residues,
    finite local densities,
    quasipolynomials,
    residue-class coverings
  }
}

\title{
Exact Formulas for Restricted Coprime Representations\\
of Even Integers with Squarefree Modulus
\texorpdfstring{$6P$}{6P}
}

\author{
Andr\'es M. Salazar\thanks{
\texttt{andresmsalazar@javerianacali.edu.co}
}\\
Department of Natural Sciences and Mathematics\\
Pontificia Universidad Javeriana Cali, Colombia
}

\date{}

\begin{document}

\maketitle

\begin{abstract}
Let $p_1,\dots,p_r\geq 5$ be distinct primes, let
$P=p_1\cdots p_r$, and put $M=6P$. For a positive integer $n$,
let $g_P(2n)$ denote the number of unordered representations
\[
2n=h+k,
\qquad
1\leq h\leq k,
\]
such that
\[
\gcd(h,M)=\gcd(k,M)=1.
\]
Using the canonical remainder operator
\[
\delta_q(x)=x-q\left\lfloor\frac{x}{q}\right\rfloor,
\]
we identify the two local residue classes excluded by each prime
$p_i$. Their possible collision is characterized exactly by
$\delta_{p_i}(n)=0$.

An inclusion--exclusion argument combined with the Chinese
remainder theorem gives an exact formula containing at most
$3^r$ terms. We prove that the associated finite correlation
has cardinality
\[
\kappa_P(n)=
\left(1+\one{\delta_3(n)=0}\right)
\prod_{i=1}^{r}
\left(p_i-2+\one{\delta_{p_i}(n)=0}\right),
\]
and obtain the affine identity
\[
g_P\bigl(2(n+6P)\bigr)-g_P(2n)=\kappa_P(n).
\]
We also establish an exact finite-density decomposition with
uniform error smaller than $3^r-1$, a half-period positivity
theorem, and a product-cutoff criterion.

Finally, we formulate the covering and paired-gap interpretations
of the problem and identify precisely the deterministic boundary
that appears when all primes not exceeding $\sqrt{2n}$ are
included. No probabilistic independence assumption is used.
\end{abstract}

\medskip

\noindent
\textbf{Keywords.}
Restricted additive representations, coprimality, Chinese
remainder theorem, canonical residues, finite local densities,
quasipolynomials, residue-class coverings.

\medskip

\noindent
\textbf{2020 Mathematics Subject Classification.}
Primary 11B75.
Secondary 11A07, 11P32.

% ============================================================
\section{Introduction}
% ============================================================

Representations of integers as sums of arithmetically restricted
summands are a classical subject of additive number theory
\cite{Nathanson1996}.
The restriction considered in this paper is finite and local.
Given
\[
P=p_1p_2\cdots p_r,
\]
where $p_1,\dots,p_r\geq5$ are distinct primes, we count
representations of $2n$ in which neither summand is divisible
by $2$, by $3$, or by any prime dividing $P$.

The one-prime case was studied in \cite{Salazar2026}. The purpose
of the present paper is to pass from the modulus $6p$ to the
squarefree modulus $6P$ while retaining an exact and explicitly
computable formula. We also verify directly that the resulting
expression reduces to the formula of \cite{Salazar2026} when
$r=1$.

There are four principal results.

First, the admissible summands are parametrized by a single
integer $z$. For every prime $p_i$, exactly two residue classes
of $z$ are potentially excluded. These two classes coincide
exactly when $p_i$ divides $n$.

Second, inclusion--exclusion and the Chinese remainder theorem
produce an exact formula with at most $3^r$ terms. For fixed $P$,
the number of arithmetic operations required by the formula does
not grow with $n$. When $r$ is regarded as a variable, the
dependence on $r$ is exponential and should not be described as
constant time.

Third, the exact local correlation has cardinality
\[
\kappa_P(n)=
\left(1+\one{\delta_3(n)=0}\right)
\prod_{i=1}^{r}
\left(p_i-2+\one{\delta_{p_i}(n)=0}\right).
\]
This quantity is simultaneously the affine increment over one
period and the cardinality defining the corresponding finite
local density. It can also be written as a finite singular factor
of Hardy--Littlewood type.

Fourth, the finite-density decomposition gives a positivity
criterion, while a separate half-period theorem yields the
product cutoff
\[
P\leq\sqrt{2n}.
\]
This product condition implies
\[
\max_{1\leq i\leq r}p_i\leq\sqrt{2n}
\]
but the converse fails when many primes are included. The open
regime is therefore the broader hypothesis in which only the
largest prime is bounded. Its extremal case is obtained by
including every prime up to $\sqrt{2n}$, and this reaches a genuine
Goldbach-type boundary.

The paper remains deterministic throughout. Statistical models
for the covering event are deliberately deferred until a sample
space, a limiting regime, and the dependence imposed by the
Chinese remainder theorem have been specified.

% ============================================================
\section{Canonical residues and the
\texorpdfstring{$6$}{6}-wheel parametrization}
% ============================================================

Throughout the paper, let
\[
P=\prod_{i=1}^{r}p_i,
\qquad
M=6P,
\]
where $r\geq1$ and $p_1,\dots,p_r\geq5$ are distinct primes.

For every positive integer $n$, define
\begin{equation}
\label{eq:g-definition}
g_P(2n)
=
\#\left\{
(h,k)\in\Z_{>0}^{\,2}
\mathrel{}\middle|\mathrel{}
h+k=2n,\ 
h\leq k,\ 
\gcd(h,M)=\gcd(k,M)=1
\right\}.
\end{equation}

Selecting the smaller summand gives the equivalent expression
\begin{equation}
\label{eq:g-short-interval}
g_P(2n)
=
\#\left\{
h\in\Z
\mathrel{}\middle|\mathrel{}
1\leq h\leq n,\ 
\gcd(h,M)=\gcd(2n-h,M)=1
\right\}.
\end{equation}
This form will be used later in the periodic and paired-gap
interpretations.

% ------------------------------------------------------------
\subsection{Canonical residues and local parameters}
% ------------------------------------------------------------

For $q\in\Z_{>0}$ and $x\in\Z$, define the canonical remainder
operator by
\begin{equation}
\label{eq:canonical-delta}
\delta_q(x)
=
x-q\left\lfloor\frac{x}{q}\right\rfloor.
\end{equation}
Thus
\[
0\leq\delta_q(x)\leq q-1.
\]
Moreover, for all $x,y\in\Z$,
\begin{equation}
\label{eq:delta-divisibility}
\delta_q(x)=\delta_q(y)
\quad\Longleftrightarrow\quad
q\mid x-y.
\end{equation}
In particular,
\[
\delta_q(x)=0
\quad\Longleftrightarrow\quad
q\mid x.
\]

The definition includes $q=1$, for which $\delta_1(x)=0$. This
convention will be useful when the Chinese-remainder modulus is
an empty product.

For every prime $p\geq5$, define
\begin{equation}
\label{eq:ab-definition}
a(p)
=
\frac{p\delta_6(p)-1}{6},
\qquad
b(p)
=
p-1-a(p).
\end{equation}

Since a prime $p\geq5$ satisfies
$\delta_6(p)\in\{1,5\}$, the parameters are explicitly given by
\[
\begin{array}{c|c|c|c}
\delta_6(p) & p & a(p) & b(p)\\
\hline
1 & 6s+1 & s & 5s\\
5 & 6s+5 & 5s+4 & s
\end{array}
\]
for a suitable nonnegative integer $s$. In particular,
$a(p),b(p)\in\{0,\dots,p-1\}$, and
\begin{align}
a(p)+b(p)
&=p-1,
\label{eq:ab-identities}
\\
6a(p)+1
&=p\delta_6(p),
\label{eq:a-identity}
\\
6b(p)+5
&=p\bigl(6-\delta_6(p)\bigr).
\label{eq:b-identity}
\end{align}

The parameters $a(p)$ and $b(p)$ identify the unique canonical
residues of $z$ for which $6z+1$ and $6z+5$, respectively, are
divisible by $p$.

\begin{lemma}
\label{lem:local-zero}
For every $z\in\Z$,
\begin{equation}
\label{eq:a-local-zero}
\delta_p(6z+1)=0
\quad\Longleftrightarrow\quad
\delta_p(z)=a(p),
\end{equation}
and
\begin{equation}
\label{eq:b-local-zero}
\delta_p(6z+5)=0
\quad\Longleftrightarrow\quad
\delta_p(z)=b(p).
\end{equation}
\end{lemma}

\begin{proof}
By \eqref{eq:a-identity}, the integer $6a(p)+1$ is divisible
by $p$. Since $\gcd(6,p)=1$, multiplication by $6$ permutes the
canonical residue classes under $\delta_p$. Consequently,
\[
p\mid6z+1
\quad\Longleftrightarrow\quad
p\mid z-a(p),
\]
which proves \eqref{eq:a-local-zero}.

The second equivalence follows in the same way from
\eqref{eq:b-identity}.
\end{proof}

% ------------------------------------------------------------
\subsection{The \texorpdfstring{$6$}{6}-wheel parametrization}
% ------------------------------------------------------------

Put
\[
e(n)=\delta_3(n)
\]
and define
\begin{equation}
\label{eq:m-definition}
m(n)
=
\begin{cases}
\dfrac{n}{3}-1,
&
e(n)=0,
\\[6pt]
\dfrac{n-1}{3},
&
e(n)=1,
\\[6pt]
\dfrac{n-5}{3},
&
e(n)=2.
\end{cases}
\end{equation}
Each branch is integral. For every positive integer $n$, one has
$m(n)\geq-1$, and equality occurs only when $n=2$. Define
\begin{equation}
\label{eq:Z-definition}
\Zcal(n)
=
\left\{
z\in\Z
\mathrel{}\middle|\mathrel{}
0\leq z\leq m(n)
\right\}.
\end{equation}
This definition automatically gives
\[
\Zcal(2)=\varnothing,
\]
because $m(2)=-1$.

For $z\in\Zcal(n)$, define $H_z(n)$ and $K_z(n)$ by
\begin{equation}
\label{eq:parametrization-table}
\begin{array}{c|c|c}
e(n) & H_z(n) & K_z(n)\\
\hline
0
&
6z+1
&
6\bigl(m(n)-z\bigr)+5
\\[4pt]
1
&
6z+1
&
6\bigl(m(n)-z\bigr)+1
\\[4pt]
2
&
6z+5
&
6\bigl(m(n)-z\bigr)+5
\end{array}
\end{equation}

In every row,
\begin{equation}
\label{eq:parametrization-sum}
H_z(n)+K_z(n)=2n.
\end{equation}
The inequalities
\[
z\geq0
\qquad\text{and}\qquad
m(n)-z\geq0
\]
ensure that both summands are positive. Conversely, positivity
of the two summands forces $0\leq z\leq m(n)$.

\begin{proposition}
\label{prop:six-wheel}
Suppose that $\delta_3(n)=0$. Then
\[
z\longmapsto\{H_z(n),K_z(n)\}
\]
is a bijection from $\Zcal(n)$ onto the unordered
representations of $2n$ whose two summands are coprime to $6$.

Suppose that $\delta_3(n)\in\{1,2\}$. Then the fibers of the same
map are precisely the orbits of the involution
\[
\tau_n(z)=m(n)-z.
\]
A fixed parameter exists exactly when $m(n)$ is even. When it
exists, it gives
\[
H_z(n)=K_z(n)=n.
\]
\end{proposition}

\begin{proof}
Every integer coprime to $6$ has canonical remainder $1$ or $5$
under $\delta_6$. Since
\[
\delta_6(2n)
=
\begin{cases}
0, & \delta_3(n)=0,\\
2, & \delta_3(n)=1,\\
4, & \delta_3(n)=2,
\end{cases}
\]
the only possible pairs of canonical residues of two summands
coprime to $6$ are
\[
(1,5)\ \text{or}\ (5,1),
\qquad
(1,1),
\qquad
(5,5),
\]
respectively.

If $\delta_3(n)=0$, exactly one summand has canonical remainder
$1$ under $\delta_6$. Writing that summand uniquely as $6z+1$
gives the first row of \eqref{eq:parametrization-table}. The
positivity conditions give precisely $z\in\Zcal(n)$. Hence every
unordered representation occurs exactly once.

If $\delta_3(n)=1$, both summands have canonical remainder $1$.
If $\delta_3(n)=2$, both have canonical remainder $5$. In either
case,
\[
H_{m(n)-z}(n)=K_z(n),
\qquad
K_{m(n)-z}(n)=H_z(n).
\]
Thus interchanging the summands corresponds exactly to replacing
$z$ by $m(n)-z$, and the fibers are the orbits of $\tau_n$.

A fixed parameter satisfies
\[
z=m(n)-z,
\]
or equivalently $2z=m(n)$. At such a parameter,
\eqref{eq:parametrization-sum} gives
\[
H_z(n)=K_z(n)=n.
\]

Finally, when $n=2$, the parameter set is empty. The only
unordered decompositions of $4$ into positive integers are
\[
4=1+3
\qquad\text{and}\qquad
4=2+2,
\]
and neither decomposition has both summands coprime to $6$.
\end{proof}

% ============================================================
\section{Local exclusions and the exact CRT formula}
% ============================================================

Proposition~\ref{prop:six-wheel} reduces the original
representation problem to the admissibility of the parameters
$z\in\Zcal(n)$. For each prime divisor $p_i$ of $P$,
divisibility of either summand excludes at most two canonical
residue classes of $z$. We now identify these classes, determine
exactly when they coincide, and combine the resulting local
conditions by inclusion--exclusion and the Chinese remainder
theorem.

For brevity, write
\[
a_i=a(p_i),
\qquad
b_i=b(p_i).
\]

For every $i\in\{1,\dots,r\}$, define the two potentially
excluded canonical residues $x_i(n)$ and $y_i(n)$ by
\begin{equation}
\label{eq:xy-definition}
\begin{array}{c|c|c}
e(n) & x_i(n) & y_i(n)\\
\hline
0
&
a_i
&
\delta_{p_i}\bigl(m(n)-b_i\bigr)
\\[4pt]
1
&
a_i
&
\delta_{p_i}\bigl(m(n)-a_i\bigr)
\\[4pt]
2
&
b_i
&
\delta_{p_i}\bigl(m(n)-b_i\bigr)
\end{array}
\end{equation}

Indeed, Lemma~\ref{lem:local-zero} and the parametrization
\eqref{eq:parametrization-table} give, for every
$z\in\Zcal(n)$,
\begin{equation}
\label{eq:local-exclusion}
p_i\mid H_z(n)K_z(n)
\quad\Longleftrightarrow\quad
\delta_{p_i}(z)\in\{x_i(n),y_i(n)\}.
\end{equation}

Thus each prime $p_i$ excludes at most two canonical residue
classes of the parameter $z$. These two classes need not always
be distinct.

\begin{lemma}[Collision lemma]
\label{lem:collision}
For every $i\in\{1,\dots,r\}$,
\begin{equation}
\label{eq:collision}
x_i(n)=y_i(n)
\quad\Longleftrightarrow\quad
\delta_{p_i}(n)=0.
\end{equation}
\end{lemma}

\begin{proof}
Suppose first that $e(n)=0$. By
\eqref{eq:xy-definition}, equality of the two canonical residues
is equivalent to
\[
\delta_{p_i}(a_i)
=
\delta_{p_i}\bigl(m(n)-b_i\bigr).
\]
Using \eqref{eq:delta-divisibility} and
$a_i+b_i=p_i-1$, this is equivalent to
\[
p_i\mid m(n)+1.
\]
Since
\[
n=3\bigl(m(n)+1\bigr)
\]
and $p_i\neq3$, the last condition is equivalent to
$\delta_{p_i}(n)=0$.

Suppose now that $e(n)=1$. Equality of the excluded residues is
equivalent to
\[
p_i\mid m(n)-2a_i.
\]
Since $n=3m(n)+1$,
\[
3\bigl(m(n)-2a_i\bigr)
=
n-\bigl(6a_i+1\bigr).
\]
By \eqref{eq:a-identity}, the quantity $6a_i+1$ is divisible by
$p_i$. Hence
\[
p_i\mid m(n)-2a_i
\quad\Longleftrightarrow\quad
p_i\mid n.
\]

Finally, suppose that $e(n)=2$. Equality is now equivalent to
\[
p_i\mid m(n)-2b_i.
\]
Because $n=3m(n)+5$,
\[
3\bigl(m(n)-2b_i\bigr)
=
n-\bigl(6b_i+5\bigr).
\]
The quantity $6b_i+5$ is divisible by $p_i$ by
\eqref{eq:b-identity}. Therefore,
\[
p_i\mid m(n)-2b_i
\quad\Longleftrightarrow\quad
p_i\mid n.
\]
This completes the proof.
\end{proof}

Define the collision indicator
\begin{equation}
\label{eq:epsilon-definition}
\varepsilon_i(n)
=
\one{\delta_{p_i}(n)=0}.
\end{equation}

The local admissibility indicator can then be written uniformly
as
\begin{equation}
\label{eq:local-indicator}
\one{
\delta_{p_i}(z)\notin\{x_i(n),y_i(n)\}
}
=
1
-
\bigl(1-\varepsilon_i(n)\bigr)
\one{\delta_{p_i}(z)=x_i(n)}
-
\one{\delta_{p_i}(z)=y_i(n)}.
\end{equation}

If $x_i(n)\neq y_i(n)$, then $\varepsilon_i(n)=0$ and the
right-hand side removes both excluded classes. If
$x_i(n)=y_i(n)$, then $\varepsilon_i(n)=1$ and the common class
is removed exactly once.

% ------------------------------------------------------------
\subsection{CRT expansion and the exact formula}
% ------------------------------------------------------------

We first record the elementary interval count used in the
expansion.

\begin{lemma}
\label{lem:L-count}
Let $m\geq-1$, let $q\geq1$, and let
$0\leq a\leq q-1$. Then
\begin{equation}
\label{eq:L-definition}
L(m,q,a)
:=
\#\left\{
z\in\Z
\mathrel{}\middle|\mathrel{}
0\leq z\leq m,\ 
\delta_q(z)=a
\right\}
=
1+\left\lfloor\frac{m-a}{q}\right\rfloor.
\end{equation}
\end{lemma}

\begin{proof}
If $a\leq m$, the integers being counted are
\[
a,\ a+q,\ a+2q,\dots
\]
up to $m$, which gives the stated formula.

If $a>m$, then
\[
-1\leq\frac{m-a}{q}<0.
\]
Consequently,
\[
1+\left\lfloor\frac{m-a}{q}\right\rfloor=0.
\]
This includes $m=-1$, for which the interval
$\{z\in\Z:0\leq z\leq m\}$ is empty.
\end{proof}

Let
\[
\boldsymbol\alpha
=
(\alpha_1,\dots,\alpha_r)
\in\{0,1,2\}^{r}
\]
and define its active index set by
\begin{equation}
\label{eq:active-indices}
I_{\boldsymbol\alpha}
=
\left\{
i\in\{1,\dots,r\}
\mathrel{}\middle|\mathrel{}
\alpha_i\neq0
\right\}.
\end{equation}

For every $i$, define the local weights
\begin{equation}
\label{eq:local-weights}
w_{i,0}(n)=1,
\qquad
w_{i,1}(n)=-(1-\varepsilon_i(n)),
\qquad
w_{i,2}(n)=-1,
\end{equation}
and put
\begin{equation}
\label{eq:global-weight}
W_{\boldsymbol\alpha}(n)
=
\prod_{i=1}^{r}w_{i,\alpha_i}(n).
\end{equation}

The modulus associated with $\boldsymbol\alpha$ is
\begin{equation}
\label{eq:Q-alpha}
Q_{\boldsymbol\alpha}
=
\prod_{i\in I_{\boldsymbol\alpha}}p_i.
\end{equation}
When $I_{\boldsymbol\alpha}=\varnothing$, the product is
understood to be $1$.

For $\ell\in\{1,2\}$, set
\begin{equation}
\label{eq:s-definition}
s_{i,1}(n)=x_i(n),
\qquad
s_{i,2}(n)=y_i(n).
\end{equation}

If $I_{\boldsymbol\alpha}=\varnothing$, define
\[
\eta_{\boldsymbol\alpha}(n)=0.
\]
If $I_{\boldsymbol\alpha}\neq\varnothing$, the Chinese remainder
theorem gives a unique
\[
\eta_{\boldsymbol\alpha}(n)
\in
\{0,\dots,Q_{\boldsymbol\alpha}-1\}
\]
satisfying
\begin{equation}
\label{eq:eta-crt}
\delta_{p_i}\bigl(\eta_{\boldsymbol\alpha}(n)\bigr)
=
s_{i,\alpha_i}(n),
\qquad
i\in I_{\boldsymbol\alpha}.
\end{equation}

For completeness, this canonical representative can be written
explicitly. For every $i\in I_{\boldsymbol\alpha}$, let
\[
Q_{\boldsymbol\alpha,i}
=
\frac{Q_{\boldsymbol\alpha}}{p_i}
\]
and let
\[
u_{\boldsymbol\alpha,i}
\in\{0,\dots,p_i-1\}
\]
be the unique canonical residue satisfying
\[
\delta_{p_i}
\left(
Q_{\boldsymbol\alpha,i}
u_{\boldsymbol\alpha,i}
\right)
=
1.
\]
Then
\begin{equation}
\label{eq:eta-explicit}
\eta_{\boldsymbol\alpha}(n)
=
\delta_{Q_{\boldsymbol\alpha}}
\left(
\sum_{i\in I_{\boldsymbol\alpha}}
s_{i,\alpha_i}(n)
Q_{\boldsymbol\alpha,i}
u_{\boldsymbol\alpha,i}
\right).
\end{equation}

Define
\begin{equation}
\label{eq:C-definition}
\Ccal_P(n)
=
\#\left\{
z\in\Zcal(n)
\mathrel{}\middle|\mathrel{}
\gcd\bigl(H_z(n)K_z(n),P\bigr)=1
\right\}
\end{equation}
and
\begin{equation}
\label{eq:D-definition}
D_P(n)
=
\one{\gcd(n,6P)=1}.
\end{equation}

\begin{theorem}[Exact CRT formula]
\label{thm:exact-formula}
For every positive integer $n$,
\begin{equation}
\label{eq:C-exact}
\Ccal_P(n)
=
\sum_{\boldsymbol\alpha\in\{0,1,2\}^{r}}
W_{\boldsymbol\alpha}(n)
L\left(
m(n),
Q_{\boldsymbol\alpha},
\eta_{\boldsymbol\alpha}(n)
\right).
\end{equation}
Moreover,
\begin{equation}
\label{eq:g-exact}
g_P(2n)
=
\begin{cases}
\Ccal_P(n),
&
\delta_3(n)=0,
\\[5pt]
\dfrac{\Ccal_P(n)+D_P(n)}{2},
&
\delta_3(n)\in\{1,2\}.
\end{cases}
\end{equation}
\end{theorem}

\begin{proof}
For $z\in\Zcal(n)$, define
\[
I_i(z)
=
\one{
\delta_{p_i}(z)\notin\{x_i(n),y_i(n)\}
}.
\]
By \eqref{eq:local-exclusion},
\[
\prod_{i=1}^{r}I_i(z)=1
\]
exactly when
\[
\gcd\bigl(H_z(n)K_z(n),P\bigr)=1.
\]
Therefore,
\begin{equation}
\label{eq:C-indicator-sum}
\Ccal_P(n)
=
\sum_{z\in\Zcal(n)}
\prod_{i=1}^{r}I_i(z).
\end{equation}

Expanding each factor by means of
\eqref{eq:local-indicator} gives
\[
\prod_{i=1}^{r}I_i(z)
=
\sum_{\boldsymbol\alpha\in\{0,1,2\}^{r}}
W_{\boldsymbol\alpha}(n)
\one{
\delta_{Q_{\boldsymbol\alpha}}(z)
=
\eta_{\boldsymbol\alpha}(n)
}.
\]
For $\boldsymbol\alpha=\boldsymbol0$, this indicator is
identically $1$ because
\[
Q_{\boldsymbol0}=1,
\qquad
\eta_{\boldsymbol0}(n)=0,
\qquad
\delta_1(z)=0.
\]

Summing the expansion over $z\in\Zcal(n)$ and applying
Lemma~\ref{lem:L-count} gives \eqref{eq:C-exact}.

Suppose that $\delta_3(n)=0$. By
Proposition~\ref{prop:six-wheel}, every admissible parameter
corresponds to exactly one unordered representation. Hence
\[
g_P(2n)=\Ccal_P(n).
\]

Suppose now that $\delta_3(n)\in\{1,2\}$. The involution
\[
z\longmapsto m(n)-z
\]
preserves admissibility because it interchanges $H_z(n)$ and
$K_z(n)$. Burnside's lemma gives
\[
g_P(2n)
=
\frac{\Ccal_P(n)+F_P(n)}{2},
\]
where $F_P(n)$ is the number of admissible fixed parameters.

A fixed parameter exists exactly when $m(n)$ is even. In the two
cases under consideration,
\[
n=3m(n)+1
\qquad\text{or}\qquad
n=3m(n)+5.
\]
Since both constant terms are odd, $m(n)$ is even exactly when
$n$ is odd. If $\gcd(n,6P)=1$, then $n$ is odd, so the fixed
parameter
\[
z=\frac{m(n)}2
\]
exists and gives $H_z(n)=K_z(n)=n$. It is admissible. Conversely,
every admissible fixed parameter gives the decomposition
$2n=n+n$ and therefore requires $\gcd(n,6P)=1$. Hence
\[
F_P(n)=D_P(n),
\]
which proves \eqref{eq:g-exact}.
\end{proof}

\begin{corollary}
\label{cor:r-one}
If $r=1$ and $P=p$, then
\begin{equation}
\label{eq:r-one-C}
\Ccal_p(n)
=
m(n)+1
-
\bigl(1-\varepsilon_1(n)\bigr)
L\bigl(m(n),p,x_1(n)\bigr)
-
L\bigl(m(n),p,y_1(n)\bigr).
\end{equation}
Together with \eqref{eq:g-exact}, this gives an equivalent
one-prime formulation of the result established in
\cite{Salazar2026}.
\end{corollary}

\begin{proof}
When $r=1$, the three possible values of $\alpha_1$ are
$0$, $1$, and $2$, with respective weights
\[
1,
\qquad
-(1-\varepsilon_1(n)),
\qquad
-1.
\]
The term corresponding to $\alpha_1=0$ is
\[
L(m(n),1,0)=m(n)+1.
\]
For $n=2$, both sides vanish because $m(2)=-1$.
\end{proof}

\begin{remark}
For fixed $P$, the inverses in \eqref{eq:eta-explicit} may be
precomputed. Formula \eqref{eq:C-exact} then contains at most
$3^r$ floor terms, and the number of arithmetic operations does
not grow with $n$. If $r$ is part of the input, the formula has
exponential dependence on $r$.
\end{remark}

% ============================================================
\section{Finite local correlations and affine periodicity}
% ============================================================

The exact formula admits a complementary interpretation at the
level of the complete modulus $M=6P$. Instead of counting
parameters in the truncated interval $\Zcal(n)$, we first count
the canonical residue classes $c$ for which both $c$ and its
reflection $c\mapsto\delta_M(2n-c)$ about $n$ are units modulo
$M$. This finite correlation determines both the local density
and the affine increment of $g_P(2n)$.

Define
\begin{equation}
\label{eq:R-definition}
\Rcal_P(n)
=
\left\{
c\in\{0,\dots,M-1\}
\mathrel{}\middle|\mathrel{}
\gcd(c,M)=1,\
\gcd\bigl(\delta_M(2n-c),M\bigr)=1
\right\}.
\end{equation}

The reflection map is
\begin{equation}
\label{eq:reflection}
\rho_n(c)
=
\delta_M(2n-c).
\end{equation}
If $c\in\Rcal_P(n)$, then $\rho_n(c)\in\Rcal_P(n)$ and
\[
\delta_M\bigl(2n-\rho_n(c)\bigr)=c.
\]
Consequently,
\[
\rho_n\bigl(\rho_n(c)\bigr)=c,
\]
so $\rho_n$ is an involution of $\Rcal_P(n)$.

Define
\begin{equation}
\label{eq:A-definition}
A_P(n)
=
\prod_{i=1}^{r}
\left(p_i-2+\varepsilon_i(n)\right)
\end{equation}
and
\begin{equation}
\label{eq:kappa-definition}
\kappa_P(n)
=
\left(1+\one{\delta_3(n)=0}\right)A_P(n).
\end{equation}
Equivalently,
\[
\kappa_P(n)
=
\left(1+\one{\delta_3(n)=0}\right)
\prod_{i=1}^{r}
\left(p_i-2+\varepsilon_i(n)\right).
\]

\begin{proposition}
\label{prop:R-cardinality}
For every positive integer $n$,
\begin{equation}
\label{eq:R-cardinality}
\#\Rcal_P(n)=\kappa_P(n).
\end{equation}
\end{proposition}

\begin{proof}
Let $q$ be a prime divisor of $M$. The condition
\[
\gcd(c,M)=1
\]
excludes the canonical residue
\[
\delta_q(c)=0.
\]
The second coprimality condition excludes
\[
\delta_q(2n-c)=0,
\]
which, by \eqref{eq:delta-divisibility}, is equivalent to
\[
\delta_q(c)=\delta_q(2n).
\]
Thus the two locally excluded residues of $c$ are
\[
0
\qquad\text{and}\qquad
\delta_q(2n).
\]

They coincide exactly when $\delta_q(2n)=0$. Hence the number of
allowed canonical residues under $\delta_q$ is
\[
q-1
\quad\text{if}\quad
\delta_q(2n)=0,
\]
and
\[
q-2
\quad\text{otherwise}.
\]

For $q=2$, the two excluded residues always coincide, giving
one allowed residue.

For $q=3$, they coincide exactly when $\delta_3(n)=0$, because
multiplication by $2$ permutes the nonzero canonical residues
under $\delta_3$. The corresponding local factor is therefore
\[
1+\one{\delta_3(n)=0}.
\]

For $q=p_i$, they coincide exactly when
$\delta_{p_i}(n)=0$. The corresponding local factor is
\[
p_i-2+\varepsilon_i(n).
\]

Since $M$ is squarefree, the Chinese remainder theorem puts the
global admissible classes in bijection with the Cartesian product
of the local admissible classes. Multiplying the local factors gives
\[
\#\Rcal_P(n)
=
\left(1+\one{\delta_3(n)=0}\right)
\prod_{i=1}^{r}
\left(p_i-2+\varepsilon_i(n)\right),
\]
which is \eqref{eq:R-cardinality}.
\end{proof}

For every prime $q\mid M$, define
\begin{equation}
\label{eq:nu-definition}
\nu_q(2n)
=
2-\one{\delta_q(2n)=0}.
\end{equation}
Thus $\nu_q(2n)$ is exactly the number of distinct local
residues excluded by the two coprimality conditions.

Define the finite singular factor
\begin{equation}
\label{eq:finite-singular-factor}
\mathfrak S_M^{\mathrm{fin}}(2n)
=
\prod_{q\mid M}
\frac{1-\nu_q(2n)/q}{(1-1/q)^2}.
\end{equation}

\begin{corollary}
\label{cor:singular-factor}
One has the exact identity
\begin{equation}
\label{eq:kappa-singular}
\frac{\kappa_P(n)}{M}
=
\left(\frac{\varphi(M)}{M}\right)^2
\mathfrak S_M^{\mathrm{fin}}(2n).
\end{equation}
Equivalently, since $M=6P$,
\[
\frac{\kappa_P(n)}{6P}
=
\left(\frac{\varphi(6P)}{6P}\right)^2
\mathfrak S_{6P}^{\mathrm{fin}}(2n).
\]
\end{corollary}

\begin{proof}
Proposition~\ref{prop:R-cardinality} gives
\[
\frac{\kappa_P(n)}{M}
=
\prod_{q\mid M}
\left(1-\frac{\nu_q(2n)}{q}\right).
\]
On the other hand,
\[
\frac{\varphi(M)}{M}
=
\prod_{q\mid M}\left(1-\frac1q\right).
\]
Factoring
\[
\left(1-\frac1q\right)^2
\]
from every local factor gives
\eqref{eq:kappa-singular}.
\end{proof}

\begin{remark}
Identity \eqref{eq:kappa-singular} is exact and finite. It does
not assert an asymptotic formula for Goldbach representations.
It identifies $\kappa_P(n)/M$ as the finite local factor
associated with the modulus $M$. The product
$\mathfrak S_M^{\mathrm{fin}}(2n)$ is the truncated analogue of
the local correction factors appearing in the
Hardy--Littlewood framework
\cite{HardyLittlewood1923,FriedlanderIwaniec2010}.
\end{remark}

% ------------------------------------------------------------
\subsection{Affine periodicity}
% ------------------------------------------------------------

\begin{theorem}[Affine periodicity]
\label{thm:affine-periodicity}
For every positive integer $n$,
\begin{equation}
\label{eq:affine-periodicity}
g_P\bigl(2(n+M)\bigr)-g_P(2n)
=
\kappa_P(n).
\end{equation}
Equivalently,
\[
g_P\bigl(2(n+6P)\bigr)-g_P(2n)
=
\kappa_P(n).
\]
\end{theorem}

\begin{proof}
Since $M=6P$, replacing $n$ by $n+M$ leaves all relevant
canonical residues unchanged:
\[
\delta_3(n+M)=\delta_3(n)
\]
and
\[
\delta_{p_i}(n+M)=\delta_{p_i}(n)
\qquad
(1\leq i\leq r).
\]
It also gives
\begin{equation}
\label{eq:m-period-increment}
m(n+M)=m(n)+2P.
\end{equation}

The local excluded classes $x_i(n)$ and $y_i(n)$ remain
unchanged. Indeed, they depend only on $a_i$, $b_i$, and the
canonical residues of $m(n)$ under the primes $p_i$, while
\[
2P
\]
is divisible by every $p_i$.

In one complete block of $P$ consecutive parameters $z$, the
number of admissible residue classes is
\[
A_P(n)
=
\prod_{i=1}^{r}
\left(p_i-2+\varepsilon_i(n)\right).
\]
This follows from the Chinese remainder theorem, since $p_i$
excludes two residue classes when $\varepsilon_i(n)=0$ and one
when $\varepsilon_i(n)=1$.

By \eqref{eq:m-period-increment}, the parameter interval for
$n+M$ contains exactly $2P$ additional consecutive integers.
It therefore contains two additional complete blocks under the
modulus $P$. Hence
\begin{equation}
\label{eq:C-period-increment}
\Ccal_P(n+M)-\Ccal_P(n)
=
2A_P(n).
\end{equation}

If $\delta_3(n)=0$, then
\[
g_P(2n)=\Ccal_P(n)
\qquad\text{and}\qquad
\kappa_P(n)=2A_P(n).
\]
Equation \eqref{eq:C-period-increment} gives the desired
increment.

If $\delta_3(n)\in\{1,2\}$, then
\[
g_P(2n)
=
\frac{\Ccal_P(n)+D_P(n)}{2}.
\]
Moreover,
\[
D_P(n+M)=D_P(n),
\]
because
\[
\gcd(n+M,M)=\gcd(n,M).
\]
It follows that
\[
g_P\bigl(2(n+M)\bigr)-g_P(2n)
=
\frac{2A_P(n)}{2}
=
A_P(n)
=
\kappa_P(n).
\]

The argument also includes $n=2$. In that case,
$\Zcal(2)=\varnothing$, while the interval associated with
$n+M$ consists of exactly $2P$ parameters.
\end{proof}

\begin{corollary}
\label{cor:quasipolynomial}
Write uniquely
\[
n=s+tM,
\qquad
1\leq s\leq M,
\qquad
t\in\Z_{\geq0}.
\]
Then
\begin{equation}
\label{eq:quasipolynomial}
g_P(2n)
=
g_P(2s)+t\kappa_P(s).
\end{equation}
Consequently, $M=6P$ is a quasiperiod of the degree-one
quasipolynomial $g_P(2n)$.
\end{corollary}

\begin{proof}
Iterate Theorem~\ref{thm:affine-periodicity} $t$ times and use
\[
\kappa_P(s+jM)=\kappa_P(s)
\qquad
(j\in\Z_{\geq0}).
\]
\end{proof}

% ------------------------------------------------------------
\subsection{The subsequence \texorpdfstring{$n=tP$}{n=tP}}
% ------------------------------------------------------------

Define the nonprincipal real Dirichlet character modulo $3$ by
\begin{equation}
\label{eq:chi3-definition}
\chi_3(a)
=
\begin{cases}
0, & \delta_3(a)=0,\\
1, & \delta_3(a)=1,\\
-1, & \delta_3(a)=2.
\end{cases}
\end{equation}

\begin{theorem}
\label{thm:ray-formula}
For every positive integer $t$,
\begin{equation}
\label{eq:ray-formula}
g_P(2tP)
=
\frac{
\left(1+\one{\delta_3(t)=0}\right)t\varphi(P)
}{6}
+
\frac{\chi_3(t)}{3}
\prod_{i=1}^{r}
\bigl(\chi_3(p_i)-1\bigr).
\end{equation}
\end{theorem}

\begin{proof}
Put $n=tP$. Formula \eqref{eq:g-short-interval} gives
\[
g_P(2tP)
=
\#\left\{
1\leq h\leq tP
\mathrel{}\middle|\mathrel{}
\gcd(h,M)=\gcd(2tP-h,M)=1
\right\}.
\]

For every prime $p_i\mid P$,
\[
p_i\mid2tP.
\]
Therefore,
\[
p_i\mid2tP-h
\quad\Longleftrightarrow\quad
p_i\mid h.
\]
The same equivalence holds for the prime $2$. Only the
restriction contributed by the prime $3$ depends on
$\delta_3(t)$.

Suppose first that $\delta_3(t)=0$. Write $t=3u$. Then
\[
2tP=uM,
\]
so
\[
\gcd(2tP-h,M)=\gcd(h,M).
\]
Consequently, $g_P(2tP)$ is the number of integers coprime to
$M$ in the interval
\[
1\leq h\leq tP=\frac{uM}{2}.
\]

The reduced residues under $\delta_M$ occur in pairs
\[
c
\qquad\text{and}\qquad
M-c.
\]
Exactly one member of each pair lies in
$[1,M/2]$. Hence
\[
\#\left\{
1\leq h\leq \frac M2
\mathrel{}\middle|\mathrel{}
\gcd(h,M)=1
\right\}
=
\frac{\varphi(M)}2
=
\varphi(P).
\]
Here $\varphi(M)=2\varphi(P)$ because $M=6P$ and
$\gcd(P,6)=1$.
Dividing the interval $[1,uM/2]$ into complete periods and,
when necessary, one half-period gives
\[
g_P(2tP)
=
\frac{u\varphi(M)}2
=
u\varphi(P)
=
\frac{t\varphi(P)}3.
\]
This is the first term of \eqref{eq:ray-formula}, while the
character term vanishes because $\chi_3(t)=0$.

Suppose now that $\delta_3(t)\neq0$. Since $P$ is coprime to
$3$, one also has $\delta_3(tP)\neq0$. The restrictions at the
primes $2$ and $3$ are equivalent to
\begin{equation}
\label{eq:ray-local-conditions}
\delta_2(h)=1,
\qquad
\delta_3(h)=\delta_3(tP).
\end{equation}
Indeed, the first condition says that $h$ is odd. Under
$\delta_3$, the two excluded residues are $0$ and
$\delta_3(2tP)$. If $\delta_3(tP)=1$, the remaining residue is
$1$, and if $\delta_3(tP)=2$, the remaining residue is $2$.

Let $\mu$ denote the M\"obius function. Applying M\"obius inversion
to the condition $\gcd(h,P)=1$ gives
\[
g_P(2tP)
=
\sum_{d\mid P}\mu(d)
\#
\left\{
1\leq h\leq tP
\mathrel{}\middle|\mathrel{}
d\mid h,\
\delta_2(h)=1,\
\delta_3(h)=\delta_3(tP)
\right\}.
\]

Write
\[
h=d\ell,
\qquad
X_d=\frac{tP}{d}.
\]
Every divisor $d$ of $P$ is coprime to $6$. In particular, $d$
is odd, so the condition under $\delta_2$ becomes
$\delta_2(\ell)=1$. Moreover, $tP=dX_d$, and multiplication by
$d$ is invertible under $\delta_3$. The second condition in
\eqref{eq:ray-local-conditions} therefore becomes
\[
\delta_3(\ell)=\delta_3(X_d).
\]

For every positive integer $X$ satisfying $\delta_3(X)\neq0$,
define
\[
N(X)
=
\#
\left\{
1\leq\ell\leq X
\mathrel{}\middle|\mathrel{}
\delta_2(\ell)=1,\
\delta_3(\ell)=\delta_3(X)
\right\}.
\]
The four possible canonical residues of $X$ under $\delta_6$
give
\[
\begin{array}{c|c|c}
\delta_6(X)
& \text{admissible residue of $\ell$ under $\delta_6$}
& N(X)\\
\hline
1 & 1 & \dfrac{X+5}{6}\\[4pt]
2 & 5 & \dfrac{X-2}{6}\\[4pt]
4 & 1 & \dfrac{X+2}{6}\\[4pt]
5 & 5 & \dfrac{X+1}{6}
\end{array}
\]
and hence
\begin{equation}
\label{eq:ray-class-count}
N(X)
=
\frac X6
+\frac12\one{\delta_2(X)=1}
+\frac{\chi_3(X)}3.
\end{equation}

Since $P/d$ is odd for every $d\mid P$, one has
\[
\delta_2(X_d)=\delta_2(t).
\]
Substitution of \eqref{eq:ray-class-count} therefore gives
\[
g_P(2tP)
=
\sum_{d\mid P}\mu(d)
\left(
\frac{tP}{6d}
+\frac12\one{\delta_2(t)=1}
+\frac{\chi_3(tP/d)}3
\right).
\]

Since $P>1$,
\[
\sum_{d\mid P}\mu(d)=0.
\]
Also,
\[
\sum_{d\mid P}\mu(d)\frac{P}{d}
=
\varphi(P),
\]
and the multiplicativity of $\chi_3$ gives
\[
\sum_{d\mid P}
\mu(d)\chi_3(tP/d)
=
\chi_3(t)
\prod_{i=1}^{r}
\bigl(\chi_3(p_i)-1\bigr).
\]
Combining these three identities proves
\eqref{eq:ray-formula}.
\end{proof}

% ============================================================
\section{Finite density, coverings, and deterministic positivity}
% ============================================================

The exact CRT formula separates naturally into a finite local
density term and a bounded fluctuation. Recall from
\eqref{eq:A-definition} and \eqref{eq:kappa-definition} that
\[
A_P(n)
=
\prod_{i=1}^{r}
\left(p_i-2+\varepsilon_i(n)\right)
\]
and
\[
\kappa_P(n)
=
\left(1+\one{\delta_3(n)=0}\right)A_P(n).
\]

For
$\boldsymbol\alpha\neq\boldsymbol0$, define
\begin{equation}
\label{eq:B-discrepancy}
\Bcal_P(n)
=
\sum_{\boldsymbol\alpha\neq\boldsymbol0}
W_{\boldsymbol\alpha}(n)
\left[
L\left(
m(n),
Q_{\boldsymbol\alpha},
\eta_{\boldsymbol\alpha}(n)
\right)
-
\frac{m(n)+1}{Q_{\boldsymbol\alpha}}
\right].
\end{equation}

\begin{proposition}
\label{prop:C-density}
For every positive integer $n$,
\begin{equation}
\label{eq:C-density}
\Ccal_P(n)
=
\frac{m(n)+1}{P}A_P(n)+\Bcal_P(n).
\end{equation}
Moreover,
\begin{equation}
\label{eq:B-bound}
\left|\Bcal_P(n)\right|
<
\prod_{i=1}^{r}
\bigl(3-\varepsilon_i(n)\bigr)-1
\leq3^r-1.
\end{equation}
\end{proposition}

\begin{proof}
From the definitions of
$W_{\boldsymbol\alpha}(n)$ and
$Q_{\boldsymbol\alpha}$,
\begin{align*}
\sum_{\boldsymbol\alpha\in\{0,1,2\}^r}
\frac{W_{\boldsymbol\alpha}(n)}
{Q_{\boldsymbol\alpha}}
&=
\prod_{i=1}^{r}
\left(
1-\frac{1-\varepsilon_i(n)}{p_i}
-\frac1{p_i}
\right)
\\
&=
\prod_{i=1}^{r}
\frac{p_i-2+\varepsilon_i(n)}{p_i}
\\
&=
\frac{A_P(n)}{P}.
\end{align*}
Separating
\[
\frac{m(n)+1}{Q_{\boldsymbol\alpha}}
\]
from each term of \eqref{eq:C-exact} gives
\eqref{eq:C-density}.

For every nonzero $\boldsymbol\alpha$, one has
$Q_{\boldsymbol\alpha}>1$. Put $N=m(n)+1$. If $N=0$, which
occurs only when $n=2$, then the discrepancy below is zero. If
$N>0$, write
\[
N=vQ_{\boldsymbol\alpha}+s,
\qquad
0\leq s<Q_{\boldsymbol\alpha}.
\]
Every canonical residue class occurs either $v$ or $v+1$ times
in any interval of $N$ consecutive integers, whereas its mean is
$v+s/Q_{\boldsymbol\alpha}$. Since
$Q_{\boldsymbol\alpha}>1$, it follows in both cases that
\[
\left|
L\left(
m(n),
Q_{\boldsymbol\alpha},
\eta_{\boldsymbol\alpha}(n)
\right)
-
\frac{m(n)+1}{Q_{\boldsymbol\alpha}}
\right|
<1.
\]

Finally,
\begin{align*}
\sum_{\boldsymbol\alpha\neq\boldsymbol0}
\left|W_{\boldsymbol\alpha}(n)\right|
&=
\prod_{i=1}^{r}
\left(
1+\left|w_{i,1}(n)\right|+\left|w_{i,2}(n)\right|
\right)-1
\\
&=
\prod_{i=1}^{r}
\bigl(3-\varepsilon_i(n)\bigr)-1.
\end{align*}
The triangle inequality now proves \eqref{eq:B-bound}.
\end{proof}

\begin{theorem}[Finite-density decomposition]
\label{thm:density-formula}
For every positive integer $n$,
\begin{equation}
\label{eq:density-formula}
g_P(2n)
=
\frac{n}{M}\kappa_P(n)+\Ecal_P(n),
\qquad
M=6P,
\end{equation}
where
\begin{equation}
\label{eq:E-branches}
\Ecal_P(n)
=
\begin{cases}
\Bcal_P(n),
&
\delta_3(n)=0,
\\[6pt]
\dfrac{\Bcal_P(n)+D_P(n)}{2}
+\dfrac{A_P(n)}{3P},
&
\delta_3(n)=1,
\\[10pt]
\dfrac{\Bcal_P(n)+D_P(n)}{2}
-\dfrac{A_P(n)}{3P},
&
\delta_3(n)=2.
\end{cases}
\end{equation}
In particular,
\begin{equation}
\label{eq:E-bound}
|\Ecal_P(n)|<3^r-1.
\end{equation}
\end{theorem}

\begin{proof}
Suppose first that $\delta_3(n)=0$. Then
\[
m(n)+1=\frac n3,
\qquad
\kappa_P(n)=2A_P(n),
\]
and $g_P(2n)=\Ccal_P(n)$. Hence
\[
g_P(2n)
=
\frac{nA_P(n)}{3P}+\Bcal_P(n)
=
\frac{n\kappa_P(n)}{6P}+\Bcal_P(n).
\]

If $\delta_3(n)=1$, then
\[
m(n)+1=\frac{n+2}{3},
\qquad
\kappa_P(n)=A_P(n).
\]
Using \eqref{eq:g-exact} and \eqref{eq:C-density}, we obtain
\begin{align*}
g_P(2n)
&=
\frac12
\left(
\frac{n+2}{3P}A_P(n)
+\Bcal_P(n)+D_P(n)
\right)
\\
&=
\frac{n\kappa_P(n)}{6P}
+
\frac{\Bcal_P(n)+D_P(n)}2
+
\frac{A_P(n)}{3P}.
\end{align*}

If $\delta_3(n)=2$, then
\[
m(n)+1=\frac{n-2}{3},
\qquad
\kappa_P(n)=A_P(n),
\]
and the same calculation gives
\[
g_P(2n)
=
\frac{n\kappa_P(n)}{6P}
+
\frac{\Bcal_P(n)+D_P(n)}2
-
\frac{A_P(n)}{3P}.
\]
This identity also includes $n=2$, for which
$m(2)+1=0$ and $D_P(2)=0$.

It remains to prove the uniform bound. When
$\delta_3(n)=0$, it follows directly from
\eqref{eq:B-bound}. In the other two cases,
\begin{align*}
|\Ecal_P(n)|
&\leq
\frac{|\Bcal_P(n)|}{2}
+\frac{D_P(n)}2
+\frac{A_P(n)}{3P}
\\
&<
\frac{3^r-1}{2}
+\frac12
+\frac13.
\end{align*}
Here
\[
0<\frac{A_P(n)}P<1
\]
because each factor of $A_P(n)$ is strictly smaller than the
corresponding prime $p_i$. Since $r\geq1$,
\[
\frac{3^r-1}{2}
+\frac12
+\frac13
<
3^r-1,
\]
which proves \eqref{eq:E-bound}.
\end{proof}

\begin{corollary}[Density criterion]
\label{cor:density-positivity}
If
\begin{equation}
\label{eq:density-criterion}
\frac{n}{M}\kappa_P(n)\geq3^r-1,
\end{equation}
then
\[
g_P(2n)>0.
\]
\end{corollary}

\begin{proof}
Theorem~\ref{thm:density-formula} gives
\[
g_P(2n)
>
\frac{n}{M}\kappa_P(n)-(3^r-1).
\]
Under \eqref{eq:density-criterion}, the right-hand side is
nonnegative. The strict inequality therefore gives
$g_P(2n)>0$.
\end{proof}

\begin{example}
The following two cases illustrate the possible signs of the
finite discrepancy:
\[
\begin{array}{c|c|c|c|c|c}
P & n & g_P(2n) & \kappa_P(n)
& \dfrac{n\kappa_P(n)}{6P} & \Ecal_P(n)\\
\hline
35 & 20 & 2 & 20 & \dfrac{40}{21} & \dfrac{2}{21}\\[6pt]
385 & 14 & 0 & 162 & \dfrac{54}{55} & -\dfrac{54}{55}
\end{array}
\]

For $P=35$ and $n=20$, the two admissible representations are
\[
40=11+29
\qquad\text{and}\qquad
40=17+23.
\]
For $P=385$ and $n=14$, there is no admissible representation,
despite the positive value of the finite local density. The
second case shows that positive local density does not by itself
force positivity on an interval that is short relative to the
period.
\end{example}

% ------------------------------------------------------------
\subsection{Coverings and paired gaps}
% ------------------------------------------------------------

For every $i\in\{1,\dots,r\}$, define
\[
X_i(n)
=
\left\{
z\in\Zcal(n)
\mathrel{}\middle|\mathrel{}
\delta_{p_i}(z)=x_i(n)
\right\}
\]
and
\[
Y_i(n)
=
\left\{
z\in\Zcal(n)
\mathrel{}\middle|\mathrel{}
\delta_{p_i}(z)=y_i(n)
\right\}.
\]

Thus $X_i(n)\cup Y_i(n)$ is the set of parameters eliminated
by the prime $p_i$.

\begin{proposition}[Covering equivalence]
\label{prop:covering}
For every positive integer $n$,
\begin{equation}
\label{eq:covering-equivalence}
g_P(2n)=0
\quad\Longleftrightarrow\quad
\Ccal_P(n)=0
\quad\Longleftrightarrow\quad
\Zcal(n)
\subseteq
\bigcup_{i=1}^{r}
\bigl(X_i(n)\cup Y_i(n)\bigr).
\end{equation}
\end{proposition}

\begin{proof}
By \eqref{eq:local-exclusion}, a parameter
$z\in\Zcal(n)$ fails to contribute to $\Ccal_P(n)$ exactly when
it belongs to at least one of the sets
$X_i(n)\cup Y_i(n)$. Therefore,
\[
\Ccal_P(n)=0
\quad\Longleftrightarrow\quad
\Zcal(n)
\subseteq
\bigcup_{i=1}^{r}
\bigl(X_i(n)\cup Y_i(n)\bigr).
\]

If $\delta_3(n)=0$, then
\[
g_P(2n)=\Ccal_P(n),
\]
so the first equivalence is immediate.

Suppose that $\delta_3(n)\in\{1,2\}$. The exact formula gives
\[
g_P(2n)
=
\frac{\Ccal_P(n)+D_P(n)}2.
\]
If $D_P(n)=1$, then $n$ is odd and coprime to $3$. Consequently,
\[
n=
\begin{cases}
6u+1, & \delta_3(n)=1,\\
6u+5, & \delta_3(n)=2,
\end{cases}
\]
for some $u\geq0$. In both cases $m(n)=2u$ is even. The central
parameter
\[
z=\frac{m(n)}2
\]
belongs to $\Zcal(n)$ and gives
\[
H_z(n)=K_z(n)=n.
\]
Because $\gcd(n,6P)=1$, this parameter is admissible. Hence
\[
D_P(n)=1
\quad\Longrightarrow\quad
\Ccal_P(n)\geq1.
\]
It follows that $\Ccal_P(n)=0$ automatically implies
$D_P(n)=0$, and the first equivalence follows.
\end{proof}

Define the periodic lift of $\Rcal_P(n)$ by
\begin{equation}
\label{eq:R-lift}
\widetilde{\Rcal}_P(n)
=
\left\{
a\in\Z
\mathrel{}\middle|\mathrel{}
\gcd(a,M)=1,\
\gcd(2n-a,M)=1
\right\}.
\end{equation}
This set contains exactly $\kappa_P(n)$ integers in every
complete interval of length $M$.

Formula \eqref{eq:g-short-interval} becomes
\begin{equation}
\label{eq:g-lift-cardinality}
g_P(2n)
=
\#
\left(
\widetilde{\Rcal}_P(n)\cap[1,n]
\right).
\end{equation}
In particular,
\begin{equation}
\label{eq:g-lift-short}
g_P(2n)>0
\quad\Longleftrightarrow\quad
\widetilde{\Rcal}_P(n)\cap[1,n]\neq\varnothing.
\end{equation}

Reflection about $n$ gives the equivalent formulation
\begin{equation}
\label{eq:g-lift-long}
g_P(2n)>0
\quad\Longleftrightarrow\quad
\widetilde{\Rcal}_P(n)\cap[1,2n-1]\neq\varnothing.
\end{equation}
Indeed, if
\[
a\in\widetilde{\Rcal}_P(n)\cap[n+1,2n-1],
\]
then
\[
2n-a\in\widetilde{\Rcal}_P(n)\cap[1,n-1].
\]

\begin{definition}
\label{def:paired-gap}
The fixed-separation paired gap $J_M(2n)$ is the least positive
integer $L$ such that every interval of $L$ consecutive
integers contains an element of
$\widetilde{\Rcal}_P(n)$. Equivalently,
\begin{equation}
\label{eq:J-definition}
J_M(2n)
=
\min
\left\{
L\in\Z_{>0}
\mathrel{}\middle|\mathrel{}
\begin{array}{c}
\text{for every }a\in\Z\text{ there exists}\\
h\in\{a+1,\dots,a+L\}
\text{ with }h\in\widetilde{\Rcal}_P(n)
\end{array}
\right\}.
\end{equation}
\end{definition}

The minimum exists because
$\widetilde{\Rcal}_P(n)$ is periodic and
$\kappa_P(n)>0$.

\begin{proposition}
\label{prop:J-bound}
For every positive integer $n$,
\begin{equation}
\label{eq:J-bound}
J_M(2n)
\leq M-\kappa_P(n)+1
\leq M-2.
\end{equation}
Moreover,
\begin{equation}
\label{eq:J-positivity}
J_M(2n)\leq2n-1
\quad\Longrightarrow\quad
g_P(2n)>0.
\end{equation}
\end{proposition}

\begin{proof}
In one circular period of length $M$, there are
$\kappa_P(n)$ admissible positions and
$M-\kappa_P(n)$ inadmissible positions. A consecutive run of
inadmissible positions therefore has length at most
$M-\kappa_P(n)$. Hence every interval of
\[
M-\kappa_P(n)+1
\]
consecutive integers contains an admissible element.

Since $r\geq1$ and every $p_i\geq5$, formula
\eqref{eq:kappa-definition} gives
\[
\kappa_P(n)\geq3.
\]
Therefore,
\[
M-\kappa_P(n)+1\leq M-2.
\]

If $J_M(2n)\leq2n-1$, the interval $[1,2n-1]$ contains an
element of $\widetilde{\Rcal}_P(n)$. The conclusion follows
from \eqref{eq:g-lift-long}.
\end{proof}

\begin{remark}
The identity
\[
\gcd(2n-h,M)=\gcd(h-2n,M)
\]
shows that membership of $h$ in $\widetilde{\Rcal}_P(n)$ is
equivalent to simultaneous coprimality of the paired progression
$(h,h-2n)$ with $M$. Thus $J_M(2n)$ is the fixed even-separation
counterpart of the paired Jacobsthal functions studied in
\cite{ZillerMorack2017,ZillerMorack2017b}. If $j_2(M)$ denotes
the corresponding global function, obtained by maximizing over
all even separations, then directly from the definitions
\[
J_M(2n)\leq j_2(M).
\]
Only the fixed-separation quantity and the elementary bound in
Proposition~\ref{prop:J-bound} are used here. No conjectural
bound for $j_2(M)$ is assumed.
\end{remark}

% ------------------------------------------------------------
\subsection{Deterministic positivity}
% ------------------------------------------------------------

\begin{theorem}[Half-period positivity]
\label{thm:half-period}
If
\begin{equation}
\label{eq:half-period-condition}
n\geq3P,
\end{equation}
then
\begin{equation}
\label{eq:half-period-positive}
g_P(2n)>0.
\end{equation}
\end{theorem}

\begin{proof}
Since $M=6P$, condition \eqref{eq:half-period-condition} is
equivalent to
\[
2n\geq M.
\]
Therefore,
\[
2n-1\geq M-1.
\]
By Proposition~\ref{prop:J-bound},
\[
J_M(2n)\leq M-2\leq2n-1.
\]
The conclusion follows from \eqref{eq:J-positivity}.
\end{proof}

\begin{proposition}[Complete-block lower bound]
\label{prop:block-bound}
For every positive integer $n$,
\begin{equation}
\label{eq:block-bound}
g_P(2n)
\geq
\left\lfloor\frac{n}{M}\right\rfloor
\kappa_P(n).
\end{equation}
Consequently, if $n\geq3P$, then
\begin{equation}
\label{eq:combined-bound}
g_P(2n)
\geq
\max
\left\{
1,\
\left\lfloor\frac{n}{M}\right\rfloor
\kappa_P(n)
\right\}.
\end{equation}
\end{proposition}

\begin{proof}
By \eqref{eq:g-lift-cardinality}, $g_P(2n)$ counts the elements
of $\widetilde{\Rcal}_P(n)$ in $[1,n]$. This interval contains
$\lfloor n/M\rfloor$ pairwise disjoint complete blocks of
length $M$, and each such block contains exactly
$\kappa_P(n)$ admissible integers. This proves
\eqref{eq:block-bound}. The second assertion follows from
Theorem~\ref{thm:half-period}.
\end{proof}

\begin{corollary}[Product cutoff]
\label{cor:product-cutoff}
If
\begin{equation}
\label{eq:product-cutoff}
P\leq\sqrt{2n},
\end{equation}
then
\[
g_P(2n)>0.
\]
\end{corollary}

\begin{proof}
Suppose first that $P\geq7$. Condition
\eqref{eq:product-cutoff} gives
\[
n\geq\frac{P^2}{2}\geq3P.
\]
The conclusion follows from Theorem~\ref{thm:half-period}.

The only remaining possibility is $P=5$. In this case,
\[
5\leq\sqrt{2n}
\]
implies $n\geq13$. For the two values not covered by the
half-period theorem, admissible representations are
\[
26=7+19
\qquad
(n=13)
\]
and
\[
28=11+17
\qquad
(n=14).
\]
For $n\geq15$, Theorem~\ref{thm:half-period} applies.
\end{proof}

\begin{remark}
The density criterion and the half-period theorem are not
comparable.

For $P=35$ and $n=106$, one has $\kappa_P(n)=15$. The
half-period theorem applies because $106\geq105$, whereas
\[
\frac{106\cdot15}{210}=\frac{53}{7}<8=3^2-1,
\]
so the density criterion does not apply.

For $P=385$ and $n=446$, one has $\kappa_P(n)=135$. In this
case
\[
\frac{446\cdot135}{2310}=\frac{2007}{77}>26=3^3-1,
\]
so the density criterion applies, whereas the half-period
condition fails because $446<1155$. Thus neither criterion
makes the other redundant.
\end{remark}

% ============================================================
\section{The prime cutoff and the deterministic boundary}
% ============================================================

For every $n\geq13$, define
\begin{equation}
\label{eq:Pi-definition}
\Pi(n)
=
\prod_{\substack{p\ \mathrm{prime}\\
5\leq p\leq\sqrt{2n}}}p.
\end{equation}
Thus $\Pi(n)$ contains every prime, other than $2$ and $3$, up
to the natural factorization threshold for integers smaller than
$2n$.

The notation $g_{\Pi(n)}(2n)$ is understood pointwise. For each
fixed $n$, the preceding theory is applied with
\[
P=\Pi(n).
\]
Since this modulus varies with $n$, the affine periodicity
theorem for a fixed $P$ is not being applied to the sequence
$g_{\Pi(n)}(2n)$.

\begin{proposition}[Monotonicity and extremality]
\label{prop:cutoff-monotonicity}
Let $P$ and $Q$ be squarefree products of primes at least $5$.
If $P\mid Q$, then, for every positive integer $n$,
\begin{equation}
\label{eq:monotonicity}
g_Q(2n)\leq g_P(2n).
\end{equation}
Consequently, if every prime factor of $P$ is at most
$\sqrt{2n}$, then
\begin{equation}
\label{eq:cutoff-extremal}
g_{\Pi(n)}(2n)\leq g_P(2n).
\end{equation}
In particular, $g_{\Pi(n)}(2n)>0$ if and only if
$g_P(2n)>0$ for every squarefree product $P$ of primes in the
interval $[5,\sqrt{2n}]$.
\end{proposition}

\begin{proof}
Every representation whose summands are coprime to $6Q$ also
has both summands coprime to $6P$. Hence the set counted by
$g_Q(2n)$ is contained in the set counted by $g_P(2n)$, which
proves \eqref{eq:monotonicity}. Under the stated cutoff on the
prime factors, one has $P\mid\Pi(n)$, and
\eqref{eq:cutoff-extremal} follows.

For the final assertion, positivity at the complete cutoff
implies positivity for every such $P$ by
\eqref{eq:cutoff-extremal}. The converse follows by taking
$P=\Pi(n)$.
\end{proof}

\begin{proposition}[Prime detection at the cutoff]
\label{prop:cutoff-prime}
Let $n\geq13$ and let
\[
2\leq u\leq2n-1.
\]
If
\[
\gcd\bigl(u,6\Pi(n)\bigr)=1,
\]
then $u$ is a prime and
\[
u>\sqrt{2n}.
\]
\end{proposition}

\begin{proof}
Suppose that $u$ is composite. Then $u$ has a prime divisor $q$
satisfying
\[
q\leq\sqrt u<\sqrt{2n}.
\]
If $q\in\{2,3\}$, then
\[
\gcd\bigl(u,6\Pi(n)\bigr)>1.
\]
If $q\geq5$, then $q$ occurs among the prime factors of
$\Pi(n)$, which gives the same contradiction. Therefore $u$
is prime.

It remains to locate this prime. If $u\leq\sqrt{2n}$, then
$u$ itself occurs among the prime factors of $6\Pi(n)$. This
again contradicts the coprimality hypothesis. Hence
$u>\sqrt{2n}$.
\end{proof}

\begin{corollary}[Cutoff characterization]
\label{cor:cutoff-equivalence}
For every $n\geq13$,
\begin{equation}
\label{eq:cutoff-equivalence}
\begin{aligned}
g_{\Pi(n)}(2n)>0
\quad\Longleftrightarrow\quad
&
\bigl(2n-1\text{ is prime}\bigr)
\\
&\text{or}
\\[-2pt]
&
\left(
\begin{array}{c}
\text{there exist primes }
q_1,q_2>\sqrt{2n}\\
\text{such that }q_1+q_2=2n
\end{array}
\right).
\end{aligned}
\end{equation}
\end{corollary}

\begin{proof}
Suppose that $g_{\Pi(n)}(2n)>0$, and let
\[
2n=h+k,
\qquad
1\leq h\leq k,
\]
be an admissible representation.

If $h=1$, then $k=2n-1$. Since $k\geq2$ and is coprime to
$6\Pi(n)$, Proposition~\ref{prop:cutoff-prime} shows that
$2n-1$ is prime.

If $h\geq2$, then both $h$ and $k$ lie between $2$ and
$2n-1$. Proposition~\ref{prop:cutoff-prime} shows that they
are primes larger than $\sqrt{2n}$. Taking
\[
q_1=h,
\qquad
q_2=k
\]
gives the second alternative.

Conversely, suppose first that $2n-1$ is prime. Since
$n\geq13$,
\[
2n-1>\sqrt{2n}>3.
\]
The prime $2n-1$ therefore divides neither $6$ nor $\Pi(n)$.
Thus
\[
2n=1+(2n-1)
\]
is an admissible representation.

Now suppose that
\[
2n=q_1+q_2,
\qquad
q_1,q_2>\sqrt{2n},
\]
with $q_1$ and $q_2$ prime. Neither prime divides
$6\Pi(n)$, because all prime divisors of this modulus are at
most $\sqrt{2n}$. After ordering the summands, this gives an
admissible representation counted by $g_{\Pi(n)}(2n)$.
\end{proof}

\begin{remark}[Precise relation with Goldbach]
If $2n-1$ is composite, then
\[
g_{\Pi(n)}(2n)>0
\]
is equivalent to the existence of a Goldbach representation
\[
2n=q_1+q_2
\]
in which both primes exceed $\sqrt{2n}$.

If $2n-1$ is prime, the admissible representation
\[
2n=1+(2n-1)
\]
makes $g_{\Pi(n)}(2n)$ positive, but it is not a Goldbach
representation because $1$ is not prime. Consequently, uniform
positivity of $g_{\Pi(n)}(2n)$ is not literally equivalent to
the full Goldbach conjecture.

Nevertheless, for those $n$ such that $2n-1$ is composite,
positivity at the complete prime cutoff is a strong
Goldbach-type assertion.
\end{remark}

\begin{remark}[Size of the cutoff product]
Let
\[
\vartheta(x)
=
\sum_{\substack{p\leq x\\p\ \mathrm{prime}}}\log p
\]
denote Chebyshev's first function. Since $n\geq13$, the primes
$2$ and $3$ lie below $\sqrt{2n}$, and therefore
\begin{equation}
\label{eq:Pi-growth}
\log\Pi(n)
=
\vartheta\bigl(\sqrt{2n}\bigr)-\log6.
\end{equation}
The prime number theorem \cite[Chapter~13]{Apostol1976} gives
\[
\vartheta(x)\sim x
\]
and hence
\begin{equation}
\label{eq:Pi-asymptotic}
\Pi(n)
=
\exp\left(
(1+o(1))\sqrt{2n}
\right)
\qquad
(n\to\infty).
\end{equation}
In particular,
\[
\frac{\Pi(n)}{n}\longrightarrow\infty.
\]
Thus the hypothesis
\[
P\leq\sqrt{2n}
\]
of Corollary~\ref{cor:product-cutoff} fails for
$P=\Pi(n)$ whenever $n$ is sufficiently large.

This explains the essential difference between the established
product condition
\[
P\leq\sqrt{2n}
\]
and the broader hypothesis
\[
\max_{1\leq i\leq r}p_i\leq\sqrt{2n}.
\]
The latter hypothesis is logically weaker, but a positivity
theorem under it would be correspondingly stronger.
For $P=\Pi(n)$, the second condition holds by construction,
while the product itself is exponentially large on the
$\sqrt n$ scale.
\end{remark}

At the complete cutoff, Proposition~\ref{prop:covering}
translates the failure of positivity into an exact covering of
the parameter interval by the local excluded classes. This is a
deterministic event. A statistical study would first have to
specify an ensemble for $n$, for the CRT-compatible vectors
$(x_i(n),y_i(n))_{1\leq i\leq r}$, or for the varying prime
cutoffs, while retaining the collision
relations, the reflection coupling, and the Chinese-remainder
constraints. Such a study is left for subsequent work. In
particular, proving that an exceptional set has density zero
would not by itself prove that the exceptional set is empty.

% ============================================================
\section{Conclusion}
% ============================================================

We have obtained an exact formula for restricted coprime
representations of even integers relative to every squarefree
modulus
\[
6P=6p_1\cdots p_r,
\]
where the primes $p_i\geq5$ are distinct. The canonical
remainder operator identifies the two local exclusions
contributed by each prime, while the collision lemma shows that
these exclusions merge exactly when
$\delta_{p_i}(n)=0$. Inclusion--exclusion and the Chinese
remainder theorem then give an exact formula with at most
$3^r$ terms.

The finite correlation $\kappa_P(n)$ determines both the affine
increment and the principal density term. The resulting
discrepancy is smaller than $3^r-1$, and the paired-gap
interpretation yields the deterministic implications
\[
n\geq3P
\quad\Longrightarrow\quad
g_P(2n)>0
\]
and
\[
P\leq\sqrt{2n}
\quad\Longrightarrow\quad
g_P(2n)>0.
\]
Among all squarefree products whose prime factors do not exceed
$\sqrt{2n}$, the complete cutoff $P=\Pi(n)$ is extremal. It lies
beyond the established product range and isolates the remaining
Goldbach-type obstruction. Its statistical analysis requires a
separate framework and is not assumed in the present work.

% ============================================================
\begin{thebibliography}{99}
% ============================================================

\bibitem{Apostol1976}
T.~M. Apostol,
\emph{Introduction to Analytic Number Theory},
Undergraduate Texts in Mathematics,
Springer, New York, 1976.

\bibitem{FriedlanderIwaniec2010}
J.~Friedlander and H.~Iwaniec,
\emph{Opera de Cribro},
American Mathematical Society Colloquium Publications,
vol.~57,
American Mathematical Society,
Providence, RI, 2010.

\bibitem{HardyLittlewood1923}
G.~H. Hardy and J.~E. Littlewood,
Some problems of ``Partitio Numerorum''.
III. On the expression of a number as a sum of primes,
\emph{Acta Mathematica}
\textbf{44} (1923), 1--70.

\bibitem{Nathanson1996}
M.~B. Nathanson,
\emph{Additive Number Theory:
The Classical Bases},
Graduate Texts in Mathematics,
vol.~164,
Springer, New York, 1996.

\bibitem{Salazar2026}
A.~M. Salazar,
Exact formulas for restricted coprime representations of even
integers,
\emph{Communications on Number Theory and Combinatorial Theory},
accepted for publication, 2026.
An earlier preprint version circulated under the title
\emph{Exact Formulas for Coprime Representations of Even
Integers Avoiding a Prime},
arXiv:2604.02386.

\bibitem{ZillerMorack2017}
M.~Ziller and J.~F. Morack,
Divisibility in paired progressions, Goldbach's conjecture,
and the infinitude of prime pairs,
arXiv:1706.00317, 2017.

\bibitem{ZillerMorack2017b}
M.~Ziller and J.~F. Morack,
A short note on the computation of the generalised Jacobsthal
function for paired progressions,
arXiv:1706.03668, 2017.

\end{thebibliography}
\end{document}
